% Part: first-order-logic
% Chapter: syntax-and-semantics
% Section: satisfaction

\documentclass[../../../include/open-logic-section]{subfiles}

\begin{document}

\olfileid{fol}{syn}{sat}

\olsection{إرضاء \article{formula} \printtoken{S}{formula}
  في \article{structure} \printtoken{S}{structure}}

\begin{explain}
يرتبط باب التعبيرات، من حدود وال!!{formula}s، بباب
ال!!{structure}s من جهتين: \emph{!!{value}} الحد، و\emph{إرضاء}
!!a{formula}. وتقريب الأولى أن !!{value} الحد !!a{element} من
!!a{structure}؛ فإن كان الحد رمزًا ثابتًا فحسب، فـ!!{value} هي الشيء
الذي عيّنته ال!!{structure} لذلك الثابت. وإن تألّف من
ال!!{function}s، حُسبت !!{value} له من !!{value}s الثوابت ومن
الدوال المعيّنة لرموز الدوال الواردة فيه. وأما !!^a{formula} فتكون
\emph{مُرضاة} في !!a{structure} متى جعل تأويل المحمولات
ال!!{formula} صادقة في مجال ال!!{structure}. وهذا أيضًا يُضبط
بالاستقراء: فتعيينات ال!!{structure} تبيّن مباشرة متى تُرضى
ال!!{formula}s الذرية؛ وأما !!{formula} المركبة، فيتعلّق إرضاؤها
بالرابط الرئيس أو المكمم، وبإرضاء ال!!{subformula}s المباشرة أو
عدم إرضائها.

غير أن في المكممات مزيد نظر؛ إذ قد يكون في ال!!{subformula}
المباشرة ل!!{formula} مكممة !!{variable} حر، وال!!{structure}s
لا تعيّن !!{value}s ال!!{variable}s. فلا بد لمعالجة ذلك من
\emph{إسنادات المتغيرات}. وحينئذ لا نعتبر الإرضاء بالنسبة إلى
!!a{structure} مجردة، بل إلى !!a{structure} ومعها إسناد
ل!!a{variable}.
\end{explain}

\begin{defn}[إسناد المتغيرات]
\emph{إسناد المتغيرات}~$s$ ل!!a{structure}~$\Struct{M}$ دالة ترسل
كل !!{variable} إلى !!a{element} من~$\Domain M$؛ أي
$s\colon \Var \to \Domain M$.
\end{defn}

\begin{explain}
في !!^a{structure} تتعيّن !!a{value} كل !!{constant}، وبإسناد
المتغيرات تتعيّن قيمة كل متغير. ونريد، فوق ذلك، أن تدل الحدود
المركبة منهما على !!{element}s من ال!!{domain}. فلذلك نعرّف
!!{value} الحدود بالاستقراء: قيمة ال!!{constant}s والمتغيرات
مأخوذة مما عيّنته ال!!{structure} أو الإسناد، وقيمة الحدود
المركبة تُحسب عوديًا بالدوال التي عيّنتها ال!!{structure}
لل!!{function}s.
\end{explain}

\begin{defn}[!!^{value} الحدود]
إذا كان $t$ حدًا من اللغة~$\Lang L$، وكانت $\Struct M$
!!{structure} لـ~$\Lang L$، وكان $s$ إسنادًا ل!!a{variable}
في~$\Struct M$، عُرّفت \emph{!!{value}}~$\Value{t}{M}[s]$ كما يأتي:
\begin{enumerate}
\item \indcase{t}{c}{$\Value{\indfrm}{M}[s] = \Assign{\indcomplex}{M}$.}
\item \indcase{t}{x}{$\Value{\indfrm}{M}[s] = s(\indcomplex)$.}
\item \indcase{t}{\Atom{f}{t_1, \ldots, t_n}}{
\[
\Value{\indfrm}{M}[s] = \Assign{f}{M}(\Value{t_1}{M}[s], \ldots,
\Value{t_n}{M}[s]).
\]}
\end{enumerate}
\end{defn}

\begin{defn}[مغاير بالنسبة إلى~$x$]
إذا كان $s$ إسنادًا ل!!a{variable} في !!a{structure}~$\Struct M$،
فكل إسناد ل!!{variable}~$s'$ في~$\Struct M$ لا يختلف عن~$s$ إلا،
على الأكثر، فيما يسنده إلى~$x$ يسمى \emph{مغايرًا بالنسبة إلى~$x$}
لـ~$s$. وإذا كان $s'$ مغايرًا بالنسبة إلى~$x$ لـ~$s$ كتبنا
$\varAssign{s'}{s}{x}$.
\end{defn}

\begin{explain}
وليس المغاير بالنسبة إلى~$x$ لإسناد~$s$ \emph{ملزمًا} بأن يجعل
قيمة~$x$ غير قيمتها الأولى؛ فكل إسناد مغاير بالنسبة
إلى~$x$ لنفسه أيضًا.
\end{explain}

\begin{defn}
  إذا كان $s$ إسنادًا ل!!a{variable} في !!a{structure}~$\Struct M$
  وكان $m \in \Domain{M}$، فإن الإسناد~$\Subst{s}{m}{x}$ هو إسناد
  المتغيرات المعرّف بـ
  \[\Subst{s}{m}{x}(y) = \begin{cases}
    m & \text{إذا كان } y \ident x\\
    s(y) & \text{فيما عدا ذلك}.
  \end{cases}\]
\end{defn}

فـ$\Subst{s}{m}{x}$ مغاير معيّن بالنسبة إلى~$x$
لـ~$s$: يضع !!{element} المجال~$m$ قيمةً لـ~$x$، ويبقي قيم
ال!!{variable}s غير~$x$ على ما عيّنه لها~$s$.

\begin{defn}[الإرضاء]
\ollabel{defn:satisfaction}
يُعرّف إرضاء !!a{formula}~$!A$ في !!a{structure}~$\Struct M$ نسبة
إلى إسناد ل!!a{variable}~$s$، ويرمز إليه
$\Sat{M}{!A}[s]$، عوديًا كما يأتي. (ونكتب
$\Sat/{M}{!A}[s]$ بمعنى «ليس $\Sat{M}{!A}[s]$».)
\begin{enumerate}
\tagitem{prvFalse}{%
  \indcase{!A}{\lfalse}{$\Sat/{M}{\indfrm}[s]$.}}{}

\tagitem{prvTrue}{%
  \indcase{!A}{\ltrue}{$\Sat{M}{\indfrm}[s]$.}}{}

\item \indcase{!A}{\Atom{R}{t_1, \dots, t_n}}{$\Sat{M}{\indfrm}[s]$
  إذا وفقط إذا كان $\langle \Value{t_1}{M}[s], \dots,
  \Value{t_n}{M}[s] \rangle \in \Assign{R}{M}$.}

\item \indcase{!A}{\eq[t_1][t_2]}{$\Sat{M}{\indfrm}[s]$ إذا وفقط
  إذا كان $\Value{t_1}{M}[s] = \Value{t_2}{M}[s]$.}

\tagitem{prvNot}{%
  \indcase{!A}{\lnot !B}{$\Sat{M}{\indfrm}[s]$ إذا وفقط إذا كان
    $\Sat/{M}{!B}[s]$.}}{}

\tagitem{prvAnd}{%
  \indcase{!A}{(!B \land !C)}{$\Sat{M}{\indfrm}[s]$ إذا وفقط إذا
    كان $\Sat{M}{!B}[s]$ و$\Sat{M}{!C}[s]$.}}{}

\tagitem{prvOr}{%
  \indcase{!A}{(!B \lor !C)}{$\Sat{M}{\indfrm}[s]$ إذا وفقط إذا
    كان $\Sat{M}{!B}[s]$ أو $\Sat{M}{!C}[s]$ (أو كلاهما).}}{}

\tagitem{prvIf}{%
  \indcase{!A}{(!B \lif !C)}{$\Sat{M}{\indfrm}[s]$ إذا وفقط إذا
    كان $\Sat/{M}{!B}[s]$ أو $\Sat{M}{!C}[s]$ (أو كلاهما).}}{}

\tagitem{prvIff}{%
  \indcase{!A}{(!B \liff !C)}{$\Sat{M}{\indfrm}[s]$ إذا وفقط إذا
    كان إما كل من $\Sat{M}{!B}[s]$ و$\Sat{M}{!C}[s]$، وإما لا
    $\Sat{M}{!B}[s]$ ولا $\Sat{M}{!C}[s]$.}}{}

\tagitem{prvAll}{%
  \indcase{!A}{\lforall[x][!B]}{$\Sat{M}{\indfrm}[s]$ إذا وفقط إذا
    كان، لكل !!{element}~$m \in \Domain M$،
    $\Sat{M}{!B}[\Subst{s}{m}{x}]$.}}{}

\tagitem{prvEx}{%
  \indcase{!A}{\lexists[x][!B]}{$\Sat{M}{\indfrm}[s]$ إذا وفقط إذا
  كان، ل!!{element} واحد على الأقل~$m \in \Domain M$،
  $\Sat{M}{!B}[\Subst{s}{m}{x}]$.}}{}
\end{enumerate}
\end{defn}

\begin{explain}
ويتبيّن موضع الحاجة إلى إسنادات المتغيرات في
\iftag{notprvEx,notprvAll}{البند الأخير}{البندين الأخيرين}.\iftag{prvAll}{
فلا يصح أن نحدّ إرضاء $\lforall[x][!B(x)]$ بقولنا: «لكل $m \in
\Domain{M}$، $\Sat{M}{!B(m)}$».}{}\iftag{prvEx}{ ولا أن نحدّ إرضاء
$\lexists[x][!B(x)]$ بقولنا: «لـ~$m \in \Domain{M}$ واحد على
الأقل، $\Sat{M}{!B(m)}$».}{} وذلك أن $m \in \Domain M$
شيء من المجال، لا رمز من اللغة؛ فلا تكون $!B(m)$~!!a{formula}،
أي إن $\Subst{!B}{m}{x}$ غير معرّفة. وليس لنا أن نفترض وجود
!!{constant}s أو حدود تسمّي كل !!{element} من~$\Struct{M}$؛
فتعريف ال!!{structure}s لا يوجب ذلك. بل إن مجموعة
ال!!{constant}s في اللغة القياسية !!{denumerable}؛ فإن كان
$\Domain{M}$ غير !!{enumerable} لم تكفِ ال!!{constant}s
لتسمية جميع الأشياء.

والخلاص من ذلك بإسنادات ال!!{variable}s، إذ نربط بها المتغيرات
مباشرة ب!!{element}s المجال. فبدلًا من أن نقول مثلًا إن
\iftag{prvEx}{$\lexists[x][!B(x)]$}{$\lforall[x][!B(x)]$} مُرضاة
في~$\Struct M$ إذا وفقط إذا
\iftag{prvEx}{كان $\Sat{M}{!B(m)}$ لواحد على الأقل من $m \in
\Domain{M}$}{كان ذلك لكل $m \in \Domain{M}$، $\Sat{M}{!B(m)}$}،
نقول: هي مُرضاة في~$\Struct M$ \emph{بالنسبة إلى}~$s$ إذا وفقط
إذا أُرضيت $!B(x)$ بالنسبة إلى~$\Subst{s}{m}{x}$
\iftag{prvEx}{لواحد على الأقل}{لكل} $m \in \Domain M$.
\end{explain}

\begin{ex}
لتكن $\Lang{L} = \{a, b, f, R\}$، حيث $a$ و$b$ !!{constant}s،
و$f$~!!{function} ثنائية المحل، و$R$~!!{predicate} ثنائي المحل. وانظر
في ال!!{structure}~$\Struct{M}$ المعرّفة بما يأتي:
\begin{enumerate}
\item $\Domain M = \{1, 2, 3, 4\}$
\item $\Assign{a}{M} = 1$
\item $\Assign{b}{M} = 2$
\item $\Assign{f}{M}(x, y) = x+y$ إذا كان $x+y \le 3$، و$= 3$ فيما عدا ذلك.
\item $\Assign{R}{M} = \{\tuple{1, 1}, \tuple{1, 2}, \tuple{2, 3}, \tuple{2, 4}\}$
\end{enumerate}
الدالة $s(x) = 1$ التي تسند $1 \in \Domain{M}$ إلى كل
!!{variable} إسناد متغيرات لـ~$\Struct{M}$.

عندئذ
\begin{align*}
\Value{f(a,b)}{M}[s] & = \Assign{f}{M}(\Value{a}{M}[s], \Value{b}{M}[s]).
\intertext{ولما كان $a$ و$b$ !!{constant}s، فإن $\Value{a}{M}[s]
  = \Assign{a}{M} = 1$ و$\Value{b}{M}[s] = \Assign{b}{M} = 2$. ومن ثم}
\Value{f(a,b)}{M}[s] & = \Assign{f}{M}(1, 2) = 1+2 = 3.
\intertext{ولحساب قيمة $f(f(a,b),a)$ يلزمنا النظر في}
\Value{f(f(a,b),a)}{M}[s] & = \Assign{f}{M}(\Value{f(a, b)}{M}[s],
\Value{a}{M}[s]) = \Assign{f}{M}(3, 1) = 3,
\intertext{لأن $3+1 > 3$. ولما كان $s(x) = 1$ و$\Value{x}{M}[s] =
  s(x)$، كان لدينا أيضًا}
\Value{f(f(a,b),x)}{M}[s] & = \Assign{f}{M}(\Value{f(a, b)}{M}[s],
\Value{x}{M}[s]) = \Assign{f}{M}(3, 1) = 3,
\end{align*}

أما ال!!{formula} الذرية~$R(t_1, t_2)$، فتُرضى متى كانت
ثنائية قيمتي حجتيها، أي $\tuple{\Value{t_1}{M}[s], \Value{t_2}{M}[s]}$،
!!a{element} من~$\Assign{R}{M}$. فمن ذلك
$\Sat{M}{R(b,f(a,b))}[s]$؛ إذ $\tuple{\Value{b}{M},
  \Value{f(a,b)}{M}} = \tuple{2, 3} \in \Assign{R}{M}$.
وأما $\Sat/{M}{R(x, f(a,b))}[s]$، فلأن
$\tuple{1, 3} \notin \Assign{R}{M}$.

وأما الصيغة غير الذرية~$!A$، فسبيل الحكم عليها أن نأخذ من
التعريف الاستقرائي البند الموافق لرابطها الرئيس. فالرابط الرئيس
في~$R(a, a) \lif (R(b, x) \lor R(x, b))$ هو~$\lif$؛ ولذلك
\begin{align*}
  & \Sat{M}{R(a, a) \lif (R(b, x) \lor R(x, b))}[s] \text{ إذا وفقط إذا }\\
  & \qquad
  \Sat/{M}{R(a,a)}[s] \text{ أو } \Sat{M}{R(b, x) \lor R(x, b)}[s]
  \intertext{ولما كان $\Sat{M}{R(a,a)}[s]$ (لأن $\tuple{1,1} \in
    \Assign{R}{M}$)، فلا يمكننا بعد تعيين الجواب، ويلزم أولًا أن
    نعرف هل $\Sat{M}{R(b, x) \lor R(x, b)}[s]$:}
  & \Sat{M}{R(b, x) \lor R(x, b)}[s] \text{ إذا وفقط إذا }\\
  & \qquad \Sat{M}{R(b,x)}[s] \text{ أو } \Sat{M}{R(x, b)}[s]
  \intertext{وهذا هو الواقع، لأن $\Sat{M}{R(x, b)}[s]$
    (إذ $\tuple{1,2} \in \Assign{R}{M}$).}
\end{align*}

قد تقدم أن المغاير بالنسبة إلى~$x$ لـ~$s$ لا يجوز أن يخالف
الإسناد~$s$ إلا في قيمة~$x$، وقد يوافقه فيها أيضًا.
فلكل !!{element} من~$\Domain{M}$ مغاير بالنسبة إلى~$x$ لـ~$s$:
\begin{align*}
  s_1 & = \Subst{s}{1}{x}, &
  s_2 & = \Subst{s}{2}{x},\\
  s_3 & = \Subst{s}{3}{x}, &
  s_4 & = \Subst{s}{4}{x}.
\end{align*}
فمثلًا، $s_2(x) = 2$ و$s_2(y) = s(y) = 1$ لجميع المتغيرات~$y$
غير~$x$. وهذه هي جميع المغايرات بالنسبة إلى~$x$ لـ~$s$ في
البنية~$\Struct{M}$، لأن $\Domain{M} = \{1, 2, 3, 4\}$. ولاحظ،
على وجه الخصوص، أن $s_1 = s$ (إذ يكون~$s$ دائمًا مغايرًا بالنسبة
إلى~$x$ لنفسه).

\iftag{prvEx}{لتعيين ما إذا كانت !!{formula} مكممة وجوديًا
~$\lexists[x][!A(x)]$ مُرضاة، يلزم أن نعيّن ما إذا كان
$\Sat{M}{!A(x)}[\Subst{s}{m}{x}]$ لواحد على الأقل من
$m \in \Domain M$. ومن ثم
  \[
  \Sat{M}{\lexists[x][(R(b,x) \lor R(x,b))]}[s],
  \]
لأن $\Sat{M}{R(b,x) \lor R(x, b)}[\Subst{s}{1}{x}]$
($\Subst{s}{3}{x}$ يصلح أيضًا). لكن
  \[
  \Sat/{M}{\lexists[x][(R(b,x) \land R(x,b))]}[s]
  \]
لأنه أيًا كان $m \in \Domain{M}$ الذي نختاره، يكون
$\Sat/{M}{R(b,x) \land R(x,b)}[\Subst{s}{m}{x}]$.}{}

\iftag{prvAll}{لتعيين ما إذا كانت !!{formula} مكممة كليًا
~$\lforall[x][!A(x)]$ مُرضاة، يلزم أن نعيّن ما إذا كان
$\Sat{M}{!A(x)}[\Subst{s}{m}{x}]$ لكل $m \in \Domain M$. ومن ثم
  \[
  \Sat{M}{\lforall[x][(R(x,a) \lif R(a,x))]}[s],
  \]
لأن $\Sat{M}{R(x,a) \lif R(a,x)}[\Subst{s}{m}{x}]$ لكل $m \in
\Domain M$. فعندما $m = 1$ يكون
$\Sat{M}{R(a,x)}[\Subst{s}{1}{x}]$، فتكون التالية صادقة؛ وعندما
$m = 2$، أو $3$، أو~$4$ يكون
$\Sat/{M}{R(x,a)}[\Subst{s}{m}{x}]$، فتكون المقدمة كاذبة. لكن
  \[
  \Sat/{M}{\lforall[x][(R(a,x) \lif R(x,a))]}[s]
  \]
لأن $\Sat/{M}{R(a,x) \lif R(x,a)}[\Subst{s}{2}{x}]$ (إذ
$\Sat{M}{R(a, x)}[\Subst{s}{2}{x}]$ و$\Sat/{M}{R(x,
  a)}[\Subst{s}{2}{x}]$).}{}

\iftag{defEx}{لتعيين ما إذا كانت !!{formula} مكممة وجوديًا
~$\lexists[x][!A(x)]$ مُرضاة، يلزم أن نعيّن ما إذا كان
$\Sat{M}{\lnot \lforall[x][\lnot !A(x)]}[s]$. فمثلًا، لدينا
  \[
  \Sat{M}{\lexists[x][(R(b,x) \lor R(x,b))]}[s].
  \]
أولًا، $\Sat{M}{R(b,x) \lor R(x, b)}[\Subst{s}{1}{x}]$
($\Subst{s}{3}{x}$ يصلح أيضًا). ومن ثم
$\Sat/{M}{\lnot(R(b,x) \lor R(x,b))}[\Subst{s}{1}{x}]$، وبالتالي
$\Sat/{M}{\lforall[x][\lnot((R(b,x) \lor R(x,b)))]}[s]$، ومنه
$\Sat{M}{\lnot\lforall[x][\lnot((R(b,x) \lor
R(x,b)))]}[s]$. وأما من الجهة الأخرى
  \[
  \Sat/{M}{\lexists[x][(R(b,x) \land R(x,b))]}[s].
  \]
وذلك لأن $\Sat{M}{\lforall[x][\lnot(R(b,x) \land
R(x,b))]}[s]$، إذ لا يوجد $m \in \Domain M$ بحيث
$\Sat{M}{R(b,x) \land R(x,b)}[\Subst{s}{m}{x}]$. ويُستدل من هذه الأمثلة على القاعدة الآتية: يكون $\Sat{M}{\lexists[x][!A(x)]}[s]$ إذا
وفقط إذا كان $\Sat{M}{!A(x)}[\Subst{s}{m}{x}]$ لواحد على الأقل
من~$m \in \Domain M$.}{}

\iftag{defAll}{لتعيين ما إذا كانت !!{formula} مكممة كليًا
~$\lforall[x][!A(x)]$ مُرضاة، يلزم أن نعيّن ما إذا كان
$\Sat{M}{\lnot\lexists[x][\lnot !A(x)]}[s]$. فمثلًا
  \[
  \Sat{M}{\lforall[x][(R(x,a) \lif R(a,x))]}[s],
  \]
أولًا، $\Sat{M}{R(x,a) \lif R(a,x)}[\Subst{s}{m}{x}]$ لكل
$m \in \Domain M$ ($\Sat{M}{R(a,x)}[\Subst{s}{1}{x}]$
و$\Sat/{M}{R(a,x)}[\Subst{s}{m}{x}]$ عندما $m = 2$، أو $3$،
أو~$4$). ومن ثم لا يوجد $m \in \Domain M$ بحيث
$\Sat{M}{\lnot(R(x,a) \lif R(a,x))}[\Subst{s}{m}{x}]$، ولذلك
$\Sat/{M}{\lexists[x][\lnot(R(x,a) \lif R(a,x))]}[s]$. وبالتالي
$\Sat{M}{\lnot\lexists[x][\lnot(R(x,a) \lif R(a,x))]}[s]$. ومن
جهة أخرى
  \[
  \Sat/{M}{\lforall[x][(R(a,x) \lif R(x,a))]}[s],
  \]
لأن $\Sat/{M}{R(a,x) \lif R(x,a)}[\Subst{s}{2}{x}]$، ومن ثم
$\Sat{M}{\lexists[x][\lnot(R(a,x) \lif R(x,a))]}[s]$. وهذه الأمثلة تدل على القاعدة الآتية: يكون
$\Sat{M}{\lforall[x][!A(x)]}[s]$ إذا وفقط إذا كان
$\Sat{M}{!A(x)}[\Subst{s}{m}{x}]$ لكل $m \in \Domain M$.}{}

ولننظر الآن في الحالة الأكثر تركيبًا:
\[
\lforall[x][(R(a,x) \lif \lexists[y][R(x,y)])].
\]
لما كان $\Sat/{M}{R(a,x)}[\Subst{s}{3}{x}]$ و
$\Sat/{M}{R(a,x)}[\Subst{s}{4}{x}]$، فلم يبق للنظر في تالية الشرطية إلا الحالتان $m = 1$
و$m = 2$. هل $\Sat{M}{\lexists[y][R(x,y)]}[\Subst{s}{1}{x}]$؟
يكون الجواب بالإيجاب إذا وُجد $n \in \Domain M$ واحد على الأقل بحيث
$\Sat{M}{R(x,y)}[\Subst{\Subst{s}{1}{x}}{n}{y}]$. ولنا أن
نأخذ $n = 1$، فيكون
$\Subst{\Subst{s}{1}{x}}{n}{y} = \Subst{s}{1}{y} = s$. ولما كان
$s(x) = 1$، و$s(y) = 1$، و$\tuple{1,1} \in \Assign{R}{M}$، فالجواب
نعم.

وأما السؤال هل
$\Sat{M}{\lexists[y][R(x,y)]}[\Subst{s}{2}{x}]$، فجوابه بالنظر في
إسنادات ال!!{variable}s~$\Subst{\Subst{s}{2}{x}}{n}{y}$. وهنا،
عندما $n = 1$، يكون هذا الإسناد~$s_2 = \Subst{s}{2}{x}$، وهذا الإسناد لا
يرضي $R(x,y)$ ($s_2(x) = 2$، و$s_2(y) = 1$،
و$\tuple{2,1}\notin \Assign{R}{M}$). وأما إذا أخذنا
$\Subst{\Subst{s}{2}{x}}{3}{y} = \Subst{s_2}{3}{y}$،
فلدينا $\Sat{M}{R(x,y)}[\Subst{s_2}{3}{y}]$ لأن
$\tuple{2,3} \in \Assign{R}{M}$، ومن ثم
$\Sat{M}{\lexists[y][R(x,y)]}[s_2]$.

وهكذا، لكل $m \in \Domain M$، إما
$\Sat/{M}{R(a,x)}[\Subst{s}{m}{x}]$ (إذا كان $m = 3$، أو~$4$)،
وإما $\Sat{M}{\lexists[y][R(x,y)]}[\Subst{s}{m}{x}]$ (إذا كان
$m = 1$، أو~$2$)، ومن ثم
\[
\Sat{M}{\lforall[x][(R(a,x) \lif \lexists[y][R(x,y)])]}[s].
\]
ومن جهة أخرى
\[
\Sat/{M}{\lexists[x][(R(a,x) \land \lforall[y][R(x,y)])]}[s].
\]
لدينا $\Sat{M}{R(a,x)}[\Subst{s}{m}{x}]$ فقط عندما $m = 1$ وعندما
$m = 2$. لكن لكل من هاتين القيمتين لـ~$m$ يوجد بدوره
$n \in \Domain M$، هو $n = 4$، بحيث
$\Sat/{M}{R(x,y)}[\Subst{\Subst{s}{m}{x}}{n}{y}]$، ومن ثم
$\Sat/{M}{\lforall[y][R(x,y)]}[\Subst{s}{m}{x}]$ عندما $m = 1$
وعندما $m = 2$. فالحاصل أنه لا يوجد $m \in \Domain M$ بحيث
$\Sat{M}{R(a,x) \land \lforall[y][R(x,y)]}[\Subst{s}{m}{x}]$.
\end{ex}

\iftag{defEx}{%
\begin{prop}\ollabel{prop:sat-ex}
$\Sat{M}{\lexists[x][!B(x)]}[s]$ إذا وفقط إذا وُجد مغاير بالنسبة
إلى~$x$، $s'$، لـ~$s$ بحيث $\Sat{M}{!B(x)}[s']$.
\end{prop}

\begin{proof}
  تمرين.
\end{proof}
}{}

\tagprob{defEx}
\begin{prob}
  أثبت \olref[fol][syn][sat]{prop:sat-ex}.
\end{prob}
\tagendprob

\iftag{defAll}{%
\begin{prop}\ollabel{prop:sat-all}
$\Sat{M}{\lforall[x][!B(x)]}[s]$ إذا وفقط إذا كان، لكل مغاير
بالنسبة إلى~$x$، $s'$، لـ~$s$، $\Sat{M}{!B(x)}[s']$.
\end{prop}

\begin{proof}
  تمرين.
\end{proof}
}{}

\tagprob{defAll}
\begin{prob}
  أثبت \olref[fol][syn][sat]{prop:sat-all}.
\end{prob}
\tagendprob

\begin{prob}
لتكن $\Lang L = \{c, f, A\}$ ذات !!{constant} واحد، و!!{function}
واحدة أحادية المحل، و!!{predicate} واحد ثنائي المحل، ولتكن
ال!!{structure}~$\Struct{M}$ معطاة بما يأتي:
\begin{enumerate}
\item $\Domain M = \{1, 2, 3\}$
\item $\Assign{c}{M} = 3$
\item $\Assign{f}{M}(1) = 2, \Assign{f}{M}(2) = 3, \Assign{f}{M}(3) = 2$
\item $\Assign{A}{M} = \{\tuple{1, 2}, \tuple{2, 3}, \tuple{3, 3}\}$
\end{enumerate}
(أ) ليكن $s(v) = 1$ لكل ال!!{variable}s~$v$. عيّن ما إذا كان
\[
\Sat{M}{\lexists[x][(A(f(z), c) \lif \lforall[y][(A(y, x) \lor A(f(y),
      x))])]}[s]
\]
وبيّن وجه الحكم.

(ب) ضع بنية أخرى وإسنادًا آخر لل!!{variable}s بحيث لا تكون
ال!!{formula} مُرضاة بهما.
\end{prob}

\end{document}
